% ----------------------------------------------------------------------------
%                                  ABSTRACT
% ----------------------------------------------------------------------------

\begin{resumo}[ABSTRACT]

    \begin{SingleSpacing}
    
       This work presents the development of a web-based system aimed at aiding behavioral management of children diagnosed with Attention Deficit Hyperactivity Disorder (ADHD) and Defiant Opposition Disorder (DOT). The central objective is to promote integrated communication between parents and teachers, in order to align strategies of accompaniment and reduce inconsistencies in the educational process. The research starts from the observation that the lack of effective communication between family and school compromises the academic and social performance of these children, in addition to generating conflicts and making it difficult to apply consistent interventions. As for the scientific methodology, the study is classified as a technological and applied research. To structure the solution, prototyping is adopted as a development method, preceded by the use of qualitative research instruments - such as the application of questionnaires, semi-structured interviews and participant observation with the school community - essential for accurate survey of requirements. From this, the project seeks to offer functionalities for the agile registration of behaviors and the monitoring of proactive pedagogical strategies. The evaluation of the tool will be conducted later through usability tests with end users, ensuring the efficiency of the interface. The expected results include improved communication between stakeholders, reduction of teacher overload and strengthening of social bonds, reinforcing the importance of technology aimed at inclusion and support in child development.
        
        \item\textbf{Keywords}: Web system. Behavioral management. ADHD. TOD. Prototyping.
    
    \end{SingleSpacing}
    
\end{resumo}
